% Part: computability
% Chapter: recursive-functions
% Section: general-recursive-functions

\documentclass[../../../include/open-logic-section]{subfiles}

\begin{document}

\olfileid{cmp}{rec}{gen}
\olsection{సామాన్య పునరావృత్త ప్రమేయాలు}

సర్వనిర్వచిత ప్రమేయాల సమితిని పొందడానికి మరొక పద్ధతి ఉంది. సహజ
సంఖ్యల ప్రతి క్రమం $\vec z$కూ $f(x,\vec z)=0$ అయ్యే $x$ ఒకటి ఉంటే,
సర్వనిర్వచిత ప్రమేయం $f(x,\vec z)$ను \emph{సక్రమం} అంటాం. మరో
విధంగా, అపరిమిత అన్వేషణను ప్రయోగించినప్పుడు సర్వనిర్వచిత ప్రమేయం
లభించే ప్రమేయాలే సక్రమ ప్రమేయాలు. ఫలితాన్ని తగ్గించకుండా అపరిమిత
అన్వేషణను సక్రమ ప్రమేయాలకే పరిమితం చేయవచ్చు:

\begin{defn}
\ollabel{defn:general-recursive}
సహజ సంఖ్యల నుంచి సహజ సంఖ్యలకు వెళ్లే వివిధ స్థానసంఖ్యల ప్రమేయాల్లో,
శూన్యం, ఉత్తరగామి, ప్రక్షేపాలను కలిగి, సంయుక్తం, ఆదిమ పునరావృత్తి,
మరియు \emph{సక్రమ} ప్రమేయాలకు ప్రయోగించిన అపరిమిత అన్వేషణ కింద
సంవృతమైన అతి చిన్న సమితిని \emph{సామాన్య పునరావృత్త ప్రమేయాల}
సమితి అంటాం.
\end{defn}

ప్రతి సామాన్య పునరావృత్త ప్రమేయం సర్వనిర్వచితమని స్పష్టం.
\olref{defn:general-recursive}, \olref[par]{defn:recursive-fn}ల మధ్య
తేడా ఏమిటంటే, రెండవ నిర్వచనంలో మధ్య దశల్లో పాక్షిక పునరావృత్త
ప్రమేయాలను వాడవచ్చు; చివరికి లభించే ప్రమేయం సర్వనిర్వచితమై ఉండటమే
ఏకైక నిబంధన. అందువల్ల ``సామాన్య'' అనే మాట చారిత్రక అవశేషమైన సరిపోని
పేరు; పైకి చూస్తే \olref{defn:general-recursive} అనేది
\olref[par]{defn:recursive-fn}కన్నా \emph{తక్కువ} సామాన్యమైనది.
అయితే అదృష్టవశాత్తూ ఆ తేడా కేవలం భ్రమ: నిర్వచనాలు వేరైనా, సామాన్య
పునరావృత్త ప్రమేయాల సమితి, పునరావృత్త ప్రమేయాల సమితి ఒకటే.

\end{document}
